% page 147 of the facsimile (image p-146 of the scan)
% block 165.1 x 233.3 mm ; left margin 36.9 mm
\pgc{15.240mm}{0.000mm}{
\l{\sou{entrechato}. (\sou{danso}).\cel{2}Salto dum qua la pedi intershokas rapide.}
\l{\cel{3}- DFIR.}
\l{}
\l{\sou{entremeso}. Disho qua venas inter la rostajo e la desero.- DEFIS.}
\l{}
\l{\sou{entr}enar. (\sou{trans., pri, por, per}). Preparar, per exerci e per}
\l{\cel{3}rejimo specala, ad ago tre energio-konsumanta. - EFS.}
\l{}
\l{\sou{entresolo}. (\sou{arkitekt.}) Lojeyo inter la ter-etajo e la etajo}
\l{\cel{3}unesma. - DEFS.}
\l{}
\l{\sou{entuziasmar}. (\sou{netrans.}) I. Sentar sua psiko exaltata. - II.}
\l{\cel{3}Ravisesar pro joyo. - DEFIRS.}
\l{}
\l{\sou{enuklear}. (\sou{trans.}) (\sou{kirurgio}). Extirpar tumoro quan, per presar,}
\l{\cel{3}onu ekirigas tra incizuro. - EF.}
\l{}
\l{\sou{enumerar}. (\sou{trans.})\cel{2}Enuncar unope omna parti di ensemblo.-EFIS.}
\l{}
\l{\sou{enuncar}. (\sou{trans}.) Adexterigar,donante a lu formo determinita,}
\l{\cel{3}to quon onu havas en sua spirito. - EFIS.}
\l{}
\l{\sou{envelopar}. (trans.) I. Cirkondar ye ulo qua kovras (lu) omna-}
\l{\cel{3}punte. - II. (\sou{Milito}). Cirkondar tale ke nula ekireyo po-}
\l{\cel{3}sibligas eskapo. - DEFI.}
\l{}
\l{\sou{enverguro}. I. Larjeso di seglo par-desvolvita. - II. Spaco quan}
\l{\cel{3}embracas la ali di ucelo (si aeroplano) kande li esas exten-}
\l{\cel{3}sita. - FS.}
\l{}
\l{\sou{envio}. Fragmenteto qua su separas de la pelo cirkum la ungli.-F.}
\l{}
\l{\sou{envidiar}. (\sou{trans.}) I. Dezirar (kozo, quan altru posedas). -}
\l{\cel{3}II. Dezirar (esar en la situeso di ulu). - EFIS.}
\l{}
\l{\sou{enzimo}. Fermento solvebla.}
\l{}
\l{\sou{eogena}. (\sou{geol.}) Dicesas pri la grupo maxim anciena ek la sulo-}
\l{\cel{3}strati terciara}
\l{}
\l{eolipilo. (\sou{tekn.}) Sfero kava, kun orifico streta, qua - plenig-}
\l{\cel{3}ite ye aquo varmigita til bolio - igas olca ekirar per sprico}
\l{\cel{3}kontinua, forta ; se instalite sur quar roti, la shoko da}
\l{\cel{3}vaporo kontre la aero extera igas la aparato avancar. -DEFIS.}
\l{}
\l{eozimo. (\sou{kemio}).Materio reda koloriziva (tetrabromofluoreceino).}
\l{\cel{3}- DEF.}
\l{}
\l{eozinofila. (\sou{citologio)}.\cel{2}Leukocito tre grosa, di qua la nukleo}
\l{\cel{3}esas polimorfa, maxim ofte duopla, mem triopla.}
\l{}
\l{epakto.(kosmol.) Nombro qua indikas quanta dii esas adjuntenda}
\l{\cel{3}a la Luno-yaro por ke lu egalesez la Suno-yaro, tale indikanta}
\l{\cel{3}la evo di Luno kande finas yaro. - DEFIRS.}
}
